\documentclass[12pt]{report}

% =====================================================================
% PREAMBLE
% =====================================================================

\usepackage[T1]{fontenc}
\usepackage[utf8]{inputenc}
\usepackage[english]{babel}

\usepackage{graphicx}
\usepackage{geometry}
\usepackage{amsmath}

% Code listings
\usepackage{listings}
\usepackage{xcolor}

\usepackage{mdframed}

% Page geometry
\geometry{
    a4paper,
    top=3cm,
    bottom=3cm,
    left=3cm,
    right=3cm
}


% =====================================================================
% CODE LISTING STYLE
% =====================================================================

\definecolor{codebackground}{RGB}{245,245,245}

\lstdefinestyle{cppstyle}{
    language=C++,
    backgroundcolor=\color{codebackground},
    basicstyle=\ttfamily\small,
    keywordstyle=\bfseries,
    commentstyle=\itshape,
    numbers=left,
    numberstyle=\tiny,
    numbersep=10pt,
    frame=single,
    framesep=6pt,
    breaklines=true,
    tabsize=4,
    showstringspaces=false,
    columns=fullflexible
}

\lstset{style=cppstyle}

% =====================================================================
% EXAMPLE BOX STYLE
% =====================================================================

\newmdenv[
    leftline=true,
    rightline=false,
    topline=false,
    bottomline=false,
    linecolor=black,
    linewidth=2pt,
    innerleftmargin=12pt,
    innerrightmargin=8pt,
    innertopmargin=8pt,
    innerbottommargin=8pt
]{examplebox}


% =====================================================================
% MAIN DOCUMENT
% =====================================================================

\begin{document}


% ---------------------------------------------------------------------
% Table of contents
% ---------------------------------------------------------------------

\clearpage
\tableofcontents
\clearpage



% =====================================================================
% INTRODUCTION
% =====================================================================

\chapter{Introduction}


Parametric polymorphism provides a language-independent abstraction, its
implementation varies considerably among programming languages and compiler
architectures.

The principal design choice concerns how generic types are represented after
compilation and when concrete type instantiations are produced.

These implementation strategies influence several aspects of the execution
model, including runtime representation, reflection capabilities, performance,
and binary size.


The following dimensions constitute the primary axes along which implementations
of parametric polymorphism are typically compared.



% ---------------------------------------------------------------------

\section{Instantiation Time}


A first distinguishing characteristic is the stage at which generic definitions
are instantiated.

Some implementations perform instantiation entirely during compilation.

In this approach, each concrete use of a generic function or data structure is
expanded into a specialized implementation before executable code is generated.

Since all specializations already exist in the final binary, runtime execution
incurs no additional generic instantiation cost.


Other implementations compile generic definitions only once.

Generic type parameters are verified during type checking and subsequently
removed or replaced by a common representation.

Consequently, no additional executable code is generated for different type
arguments, and all instantiations share the same compiled implementation.



% ---------------------------------------------------------------------

\section{Runtime Representation}


Implementations also differ in how generic types are represented after
compilation.

\emph{Type erasure} removes generic type parameters, producing a homogeneous
representation where different instantiations share the same executable code.

The runtime system therefore treats different generic instantiations as
equivalent representations.


Conversely, \emph{reified generics} preserve concrete type information by
generating distinct specialized code for each instantiated type.

Each generic instantiation may therefore possess its own runtime
representation, memory layout, and executable instructions.



% ---------------------------------------------------------------------

\section{Reflection and Runtime Type Information}


The availability of generic type information at runtime depends directly on the
chosen representation strategy.


When type erasure is employed, generic parameters are generally unavailable after
compilation.

Reflection mechanisms therefore provide limited access to instantiated generic
types because most type information has been discarded.


Conversely, implementations based on reified generics preserve concrete type
information throughout compilation.

\emph{Runtime reflection or runtime type identification (RTTI)} may therefore
recover the actual instantiated types, enabling operations such as generic type
inspection, specialized serialization, and runtime metadata analysis.



% ---------------------------------------------------------------------

\section{Type Safety}


Both implementation strategies preserve static type safety, although they
achieve this objective differently.


In implementations based on type erasure, generic correctness is verified
entirely during compilation.

Although type parameters disappear from the generated executable,
compiler-inserted casts preserve the guarantees established by the static type
system.


Implementations based on reified generics retain complete type information
throughout compilation and generate type-specific implementations directly.

Since each specialization corresponds to a concrete type, correctness is
preserved without requiring erased representations or implicit runtime casts.



% ---------------------------------------------------------------------

\section{Performance}


The implementation strategy significantly influences both compilation and
execution performance.


Homogeneous implementations generally produce smaller binaries and reduce
compilation time because generic definitions are compiled only once.

However, the shared representation may require additional indirections,
boxing, or runtime casts, limiting optimization opportunities.


Heterogeneous implementations incur greater compilation cost and larger
executables because every instantiated type generates a distinct implementation.

Nevertheless, the compiler can aggressively optimize each specialization,
exploiting concrete type information for improved instruction selection, memory
layout, inlining, and register allocation.


Consequently, runtime performance often approaches that of manually specialized
code.



% ---------------------------------------------------------------------

\section{Summary}


The implementation of parametric polymorphism is primarily characterized by five
orthogonal dimensions:

\begin{enumerate}
    \item the stage at which generic instantiation occurs;
    \item the runtime representation of generic types;
    \item the availability of runtime type information and reflection;
    \item the trade-off between compilation cost, binary size, and execution
          performance;
    \item the mechanisms through which static type safety is preserved.
\end{enumerate}


These dimensions form the basis for comparing homogeneous implementations based
on type erasure with heterogeneous implementations based on reification or
monomorphization.




\chapter{Instantiation Time}


\section{Instantiation Time in C++ Templates}


C++ implements parametric polymorphism through the \emph{template} mechanism,
which adopts a purely compile-time instantiation model.

Unlike languages that compile a generic definition only once, a C++ template is
not itself an executable entity. Instead, it represents a parameterized program definition that serves as a
blueprint from which the compiler generates independent concrete
implementations.

% ---------------------------------------------------------------------

From the perspective of instantiation time, the defining characteristic of C++
is that executable code is generated only after the compiler knows the complete
set of template arguments.

Consequently, all template instantiations are resolved entirely during static
compilation, and no template instantiation mechanism exists during program
execution.


The complete instantiation process can be understood by following the template
through the compilation pipeline.

% ---------------------------------------------------------------------

\subsection{Template Definition}

When the compiler encounters a template definition, it first performs the
ordinary front-end compilation stages, including lexical analysis, parsing,
and semantic analysis that can be performed without knowledge of future
template arguments.


For example:

\begin{lstlisting}
template<typename T>
T get_max(T a, T b)
{
    return (a > b) ? a : b;
}
\end{lstlisting}


The compiler parses the template syntax, builds the corresponding Abstract
Syntax Tree (AST), checks syntactic correctness, and stores an internal
parameterized representation.

At this stage, the compiler does not generate executable machine code.

The reason is straightforward: the compiler still ignores the concrete meaning
of the template parameter \texttt{T}, and therefore cannot determine the exact
operations, object layout, calling convention, or generated instructions
required by the final implementation.

Consequently, the template definition remains a generic blueprint waiting for
future instantiation.

% ---------------------------------------------------------------------

\subsection{Instantiation Trigger}

The instantiation process begins only when the compiler encounters a concrete
use of the template, either through an explicit specification of the template
argument or through implicit type deduction.

\begin{examplebox}

\textbf{Example.}

Consider the following template invocations:

\begin{lstlisting}
int i = get_max(3, 7);
double d = get_max<double>(2.5, 4.8);
\end{lstlisting}


The first invocation relies on template argument deduction: 
\[
T = int
\]


whereas the second invocation explicitly specifies the template argument:

\[
T = double
\]


Therefore, the compiler generates two different template instantiations:
\texttt{get\_max<int>} and \texttt{get\_max<double>}

\end{examplebox}
Each unique set of template arguments constitutes a distinct template
instantiation.
Instantiation may occur through:

\begin{itemize}
    \item \texttt{implicit instantiation}, automatically performed whenever a
          specialization is required by the program;

    \item \texttt{explicit instantiation}, requested directly by the programmer through
          an explicit instantiation declaration or definition;

    \item \texttt{explicit specialization}, where an entirely different implementation
          is provided for a specific combination of template arguments.
\end{itemize}

Regardless of how instantiation is initiated, executable code generation begins
only after the complete set of template arguments is known.

% ---------------------------------------------------------------------

\subsection{Template Argument Substitution}


Once instantiation has been triggered, the compiler performs template argument
substitution.


During this phase, every formal template parameter is replaced by its
corresponding concrete argument.

\begin{examplebox}
\textbf{From previous example.}

\[
T \rightarrow int
\]


produces the specialization:

\begin{lstlisting}
int get_max(int a, int b)
{
    return (a > b) ? a : b;
}
\end{lstlisting}


whereas:

\[
T \rightarrow double
\]


produces:

\begin{lstlisting}
double get_max<double>(double a, double b)
{
    return (a > b) ? a : b;
}
\end{lstlisting}
\end{examplebox}

Conceptually, the compiler now treats each specialization as if it had been
written explicitly by the programmer.

From this point onward, the specialization behaves as an ordinary non-template
function.

% ---------------------------------------------------------------------

\subsection{Semantic Analysis of the Specialization}

Instantiation does not immediately produce executable code.

Instead, each generated specialization undergoes the normal semantic analysis
performed for ordinary C++ program entities.

Depending on the specialization, the compiler performs activities such as: type checking, overload resolution, name lookup, constraint verification for constrained templates,
constant expression evaluation, instantiation of nested templates and verification of language rules required by the selected template arguments.

This step is fundamental because many semantic properties cannot be verified
before template parameters have been substituted.


Consequently, a template definition that is syntactically correct may still
produce compilation errors if one of its concrete instantiations violates the
requirements imposed by the template body.

% ---------------------------------------------------------------------

\subsection{Intermediate Representation and Code Generation}

Once semantic analysis has successfully completed, the specialization enters
the compiler back-end.

The compiler generates an Intermediate Representation (IR) specifically for the
instantiated specialization.

Subsequently, the specialization undergoes the standard optimization pipeline,
which may include: constant propagation, dead-code elimination, function inlining, loop optimizations, register allocation,
instruction selection and target-specific optimizations.

Finally, the compiler emits native machine code corresponding exclusively to
that specialization.

Importantly, this process is repeated independently for every distinct template
instantiation.

Consequently, each specialization owns an independent executable
implementation. Infact, the compiler does not produce a single shared executable representation capable
of serving multiple template arguments.

% ---------------------------------------------------------------------

\subsection{Completion of the Instantiation Process}

The instantiation process terminates entirely before linking.

By the time object files are generated, every template specialization required
by the translation unit already exists as ordinary compiled code.

After linking, the final executable therefore contains every specialization
required by the program.

From the perspective of executable code generation, the lifecycle of a C++
template can be summarized as follows:

\begin{enumerate}
    \item The compiler parses the template definition.
    
    \item The template is stored as a parameterized internal representation.

    \item A concrete template use triggers instantiation.

    \item Template parameters are substituted with concrete arguments.

    \item The generated specialization undergoes complete semantic analysis.

    \item The compiler produces a specialization-specific intermediate
          representation.

    \item Optimization passes are applied.

    \item Native machine code is emitted.

    \item The specialization is included in the final executable before program
          execution begins.
\end{enumerate}

Therefore, the defining characteristic of C++ parametric polymorphism is that
the entire instantiation process is completed during static compilation.

When execution begins, no template instantiation mechanism remains active.

The runtime simply executes machine code that has already been generated for
every required specialization during compilation.

% =====================================================================
% FUTURE SECTIONS
% =====================================================================

\section{Analysis of the Algorithms}


You can use sections to divide your chapters into smaller, readable chunks.


\end{document}
